\documentclass[11pt,a4paper]{article}

% ===== XeLaTeX + Korean =====
\usepackage{kotex}
\usepackage{fontspec}
\setmainfont{Times New Roman}
\setmainhangulfont{AppleMyungjo}
\setsanshangulfont{Apple SD Gothic Neo}

% ===== Layout =====
\usepackage[a4paper,top=18mm,bottom=18mm,left=20mm,right=20mm]{geometry}
\usepackage{fancyhdr}
\setlength{\headheight}{24pt}
\usepackage{enumitem}
\usepackage{mdframed}
\usepackage{array}
\usepackage{booktabs}

\setlength{\parindent}{1em}
\linespread{1.18}

% ===== Header / Footer =====
\pagestyle{fancy}
\fancyhf{}
\renewcommand{\headrulewidth}{0.6pt}
\fancyhead[L]{\sffamily\bfseries 2026 고1 영어 변형 --- 지문 33}
\fancyhead[R]{\sffamily 후각 묘사의 어려움}
\renewcommand{\footrulewidth}{0pt}
\fancyfoot[C]{\framebox[16mm][c]{\thepage}}

% ===== helpers =====
\newcommand{\qno}[1]{\par\vspace{8pt}\noindent{\sffamily\bfseries #1.}\ }
\newcommand{\instr}[1]{{\sffamily #1}\par\vspace{3pt}}
\newlist{choices}{enumerate}{1}
\setlist[choices]{label=,leftmargin=1.5em,itemsep=2pt,topsep=3pt,parsep=0pt}

\newmdenv[linewidth=0.6pt,roundcorner=3pt,innertopmargin=7pt,innerbottommargin=7pt,
          innerleftmargin=9pt,innerrightmargin=9pt,skipabove=5pt,skipbelow=5pt]{pbox}
\newcommand{\nemo}[1]{\fbox{\,#1\,}}

\begin{document}

\begin{center}
{\fontsize{15}{17}\selectfont\sffamily\bfseries [지문 33] 다음 글을 읽고 물음에 답하시오.}
\end{center}
\vspace{4pt}

% ===== 공통 지문 =====
\begin{pbox}
Talking about what something smells like is the most basic way of talking about our olfactory$^{*}$ experience, and even this most basic way of talking about smells is \textcircled{\scriptsize 1}\,\underline{challenging}. There are no words in the English language to describe smells in the same way in which ``blue'' or ``green'' describe colors. Instead, to talk about how something smells, we talk about \rule{3.5cm}{0.4pt}. Things smell ``flowery,'' ``fruity,'' or ``fishy.'' Furthermore, even the most \textcircled{\scriptsize 3}\,\underline{unfamiliar} odors are difficult to \textcircled{\scriptsize 2}\,\underline{identify} when they are not experienced in their usual context. In one experiment the majority of participants was \textcircled{\scriptsize 4}\,\underline{unable} to name very common odors like beer, urine$^{**}$, roses, or motor oil. Obviously, even those who couldn't name any of these odors would drink the beer but not the urine or the motor oil. This is how evolution has \textcircled{\scriptsize 5}\,\underline{shaped} our brain: We respond to odors correctly in many different ways, but we are not well-equipped to talk about them.

\vspace{4pt}
\noindent{\small $^{*}$ olfactory: 후각의\quad $^{**}$ urine: 소변}
\end{pbox}

% ===================== Q1 (빈칸추론 3점) =====================
\qno{1}\instr{다음 빈칸에 들어갈 말로 가장 적절한 것은? [3점]}
\begin{choices}
\item ① how pleasant the odor is
\item ② how well we can identify it
\item ③ what produces the smell
\item ④ the emotional reaction it triggers
\item ⑤ how familiar the smell is to us
\end{choices}

% ===================== Q2 (어휘 부적절) =====================
\qno{2}\instr{밑줄 친 ①$\sim$⑤ 중에서 문맥상 낱말의 쓰임이 적절하지 않은 것은?}
\begin{choices}
\item ①\quad ②\quad ③\quad ④\quad ⑤
\end{choices}

% ===================== Q3 (서술형 단답 한글) =====================
\qno{3}\instr{[서술형] 윗글의 실험(experiment)에서 발견된 결과를 한 문장(30자 이내)으로 우리말로 쓰시오.}
\par\vspace{4pt}\noindent
$\Rightarrow$ \rule{10cm}{0.4pt}

% ===================== Q4 =====================
\qno{4}\instr{다음은 윗글을 요약한 것이다. 빈칸 (A), (B)에 들어갈 말로 가장 적절한 것은?}
\noindent\hspace{1em}\textit{Although humans can respond to odors (A)\,\rule{2.5cm}{0.4pt}\, in practical situations, they often have difficulty (B)\,\rule{2.5cm}{0.4pt}\, them in language.}
\begin{center}
\renewcommand{\arraystretch}{1.2}
\begin{tabular}{c c c}
 & (A) & (B) \\
① & correctly & describing \\
② & randomly & avoiding \\
③ & emotionally & producing \\
④ & visually & measuring \\
⑤ & slowly & preserving \\
\end{tabular}
\end{center}

\end{document}
